% GENERATED FILE. Edit canonical lessons, then run npm run generate:books.
% canonical-source-hash: fnv1a64:8c761747ed16684f
% canonical-lessons: TE-C419-kakapote, TE-C419-batti, TE-C419-minahayimpu, TE-C419-emainappatiki, TE-C419-sambhavincu

\chapter{If Not, Then What}
\label{ch:te-if-not-then-what}
\hlchaptermodality{419}

\begin{chapteropening}
\textbf{By the end of this chapter:} \emph{I can say otherwise, depending on, and whatever happens.}

Complete the last lesson of chapter 419: I can say otherwise, depending on, and whatever happens.
\end{chapteropening}

\section[\texorpdfstring{\te{కాకపోతే} (kākapōtē)}{కాకపోతే (kākapōtē)}]{\te{కాకపోతే} (kākapōtē) — if not; otherwise}

\label{lesson:TE-C419-kakapote}



\begin{quote}
Before the new one: say the Telugu for unavoidable, then the Telugu for if possible.
\end{quote}

\subsection*{You'll want to know: \te{కాకపోతే}}
\textbf{\te{కాకపోతే}} — \emph{kākapōtē} — \textquotedblleft{}if not; otherwise\textquotedblright{}.

\begin{quote}
\te{బస్సు} \te{దొరికితే} \te{వస్తాను}, \te{కాకపోతే} \te{రేపు} \te{వస్తాను}.
\end{quote}

\begin{quote}
\emph{bassu dorikitē vastānu, kākapōtē rēpu vastānu.}
\end{quote}

\begin{quote}
If I get a bus I will come; if not, I will come tomorrow.
\end{quote}

A cousin of \te{లేకపోతే}: if it does not turn out so.

\begin{cousinweb}
\textbf{\te{కాకపోతే}} (\emph{kākapōtē}) is native Dravidian: \emph{kā-}, a stem of \emph{avu}, \textquotedblleft{}to be\textquotedblright{} (DEDR 333), and \emph{pōtē}, \textquotedblleft{}if it goes\textquotedblright{}, a form of \emph{pōvu}, \textquotedblleft{}to go\textquotedblright{} (DEDR 4572).
\end{cousinweb}

\begin{grammarlens}[title={What you've built}]
One more word for saying what happens if.
\end{grammarlens}

\subsection*{Your turn}
\begin{itemize}
\raggedright
  \item \emph{Say it:} \emph{kākapōtē}
  \item \emph{Say it:} \emph{kākapōtē}, once more
  \item \emph{Say it:} \emph{bassu dorikitē vastānu, kākapōtē rēpu vastānu.} — If I get a bus I will come; if not, I will come tomorrow.
  \item \emph{Recall:} say the Telugu for unavoidable, then the Telugu for if possible, then say \emph{kākapōtē} again
\end{itemize}

\begin{tcolorbox}[breakable,colback=teal!4,colframe=teal!35!black,title={Before you move on}]
What does \te{కాకపోతే} mean? (\textquotedblleft{}If not; otherwise\textquotedblright{}.) Say it once more.
\end{tcolorbox}
\section[\texorpdfstring{\te{బట్టి} (baṭṭi)}{బట్టి (baṭṭi)}]{\te{బట్టి} (baṭṭi) — depending on; going by}

\label{lesson:TE-C419-batti}



\begin{quote}
Before the new one: say the Telugu for if possible, then the Telugu for if not.
\end{quote}

\subsection*{You'll want to know: \te{బట్టి}}
\textbf{\te{బట్టి}} — \emph{baṭṭi} — \textquotedblleft{}depending on; going by\textquotedblright{}.

\begin{quote}
\te{ధరను} \te{బట్టి} \te{నేను} \te{కొంటాను}.
\end{quote}

\begin{quote}
\emph{dharanu baṭṭi nēnu koṇṭānu.}
\end{quote}

\begin{quote}
I will buy it depending on the price.
\end{quote}

The thing depended on takes -\te{ని}: \te{ధరను} \te{బట్టి}.

\begin{cousinweb}
\textbf{\te{బట్టి}} (\emph{baṭṭi}) is native Dravidian, usually taken as a form of \emph{paṭṭu}, \textquotedblleft{}to hold, to take hold of\textquotedblright{} (DEDR 4034): holding to a thing, going by it.
\end{cousinweb}

\begin{grammarlens}[title={What you've built}]
One more word for saying what happens if.
\end{grammarlens}

\subsection*{Your turn}
\begin{itemize}
\raggedright
  \item \emph{Say it:} \emph{baṭṭi}
  \item \emph{Say it:} \emph{baṭṭi}, once more
  \item \emph{Say it:} \emph{dharanu baṭṭi nēnu koṇṭānu.} — I will buy it depending on the price.
  \item \emph{Recall:} say the Telugu for if possible, then the Telugu for if not, then say \emph{baṭṭi} again
\end{itemize}

\begin{tcolorbox}[breakable,colback=teal!4,colframe=teal!35!black,title={Before you move on}]
What does \te{బట్టి} mean? (\textquotedblleft{}Depending on; going by\textquotedblright{}.) Say it once more.
\end{tcolorbox}
\section[\texorpdfstring{\te{మినహాయింపు} (minahāyimpu)}{మినహాయింపు (minahāyimpu)}]{\te{మినహాయింపు} (minahāyimpu) — an exemption}

\label{lesson:TE-C419-minahayimpu}



\begin{quote}
Before the new one: say the Telugu for if not, then the Telugu for depending on.
\end{quote}

\subsection*{You'll want to know: \te{మినహాయింపు}}
\textbf{\te{మినహాయింపు}} — \emph{minahāyimpu} — \textquotedblleft{}an exemption\textquotedblright{}.

\begin{quote}
\te{విద్యార్థులకు} \te{ఫీజులో} \te{మినహాయింపు} \te{ఇచ్చారు}.
\end{quote}

\begin{quote}
\emph{vidyārthulaku phījulō minahāyimpu iccāru.}
\end{quote}

\begin{quote}
Students were given an exemption from the fee.
\end{quote}

The noun of \te{మినహాయించు}, to exempt, which is built on \te{మినహా}, except: -\te{ఇంపు} makes a noun of an -\te{ఇంచు} verb.

\begin{cousinweb}
\textbf{\te{మినహాయింపు}} (\emph{minahāyimpu}) is built on \te{మినహా}, \textquotedblleft{}except\textquotedblright{}, which is usually traced to Hindustani \emph{minhā}, of Arabic origin, with -\te{ఇంపు}, which makes a noun of the verb \te{మినహాయించు}, to exempt.
\end{cousinweb}

\begin{grammarlens}[title={What you've built}]
One more word for saying what happens if.
\end{grammarlens}

\subsection*{Your turn}
\begin{itemize}
\raggedright
  \item \emph{Say it:} \emph{minahāyimpu}
  \item \emph{Say it:} \emph{minahāyimpu}, once more
  \item \emph{Say it:} \emph{vidyārthulaku phījulō minahāyimpu iccāru.} — Students were given an exemption from the fee.
  \item \emph{Recall:} say the Telugu for if not, then the Telugu for depending on, then say \emph{minahāyimpu} again
\end{itemize}

\begin{tcolorbox}[breakable,colback=teal!4,colframe=teal!35!black,title={Before you move on}]
What does \te{మినహాయింపు} mean? (\textquotedblleft{}An exemption\textquotedblright{}.) Say it once more.
\end{tcolorbox}
\section[\texorpdfstring{\te{ఏమైనప్పటికీ} (ēmainappaṭikī)}{ఏమైనప్పటికీ (ēmainappaṭikī)}]{\te{ఏమైనప్పటికీ} (ēmainappaṭikī) — in any case; whatever happens}

\label{lesson:TE-C419-emainappatiki}



\begin{quote}
Before the new one: say the Telugu for depending on, then the Telugu for an exemption.
\end{quote}

\subsection*{You'll want to know: \te{ఏమైనప్పటికీ}}
\textbf{\te{ఏమైనప్పటికీ}} — \emph{ēmainappaṭikī} — \textquotedblleft{}in any case; whatever happens\textquotedblright{}.

\begin{quote}
\te{ఏమైనప్పటికీ} \te{నేను} \te{పరీక్ష} \te{రాస్తాను}.
\end{quote}

\begin{quote}
\emph{ēmainappaṭikī nēnu parīkṣa rāstānu.}
\end{quote}

\begin{quote}
Whatever happens, I will sit the exam.
\end{quote}

\te{ఏమి}, what, + \te{అయినప్పటికీ}, even if it is. Close to \te{ఏదేమైనా}, which you know.

\begin{cousinweb}
\textbf{\te{ఏమైనప్పటికీ}} (\emph{ēmainappaṭikī}) is native Dravidian: \emph{ēmi}, \textquotedblleft{}what\textquotedblright{} (DEDR 5151), and \emph{ayinappaṭikī}, \textquotedblleft{}even though it is\textquotedblright{}, built on \emph{avu}, \textquotedblleft{}to be\textquotedblright{} (DEDR 333).
\end{cousinweb}

\begin{grammarlens}[title={What you've built}]
One more word for saying what happens if.
\end{grammarlens}

\subsection*{Your turn}
\begin{itemize}
\raggedright
  \item \emph{Say it:} \emph{ēmainappaṭikī}
  \item \emph{Say it:} \emph{ēmainappaṭikī}, once more
  \item \emph{Say it:} \emph{ēmainappaṭikī nēnu parīkṣa rāstānu.} — Whatever happens, I will sit the exam.
  \item \emph{Recall:} say the Telugu for depending on, then the Telugu for an exemption, then say \emph{ēmainappaṭikī} again
\end{itemize}

\begin{tcolorbox}[breakable,colback=teal!4,colframe=teal!35!black,title={Before you move on}]
What does \te{ఏమైనప్పటికీ} mean? (\textquotedblleft{}In any case; whatever happens\textquotedblright{}.) Say it once more.
\end{tcolorbox}
\section[\texorpdfstring{\te{సంభవించు} (sambhaviñcu)}{సంభవించు (sambhaviñcu)}]{\te{సంభవించు} (sambhaviñcu) — to occur; to take place}

\label{lesson:TE-C419-sambhavincu}



\begin{quote}
Before the new one: say the Telugu for an exemption, then the Telugu for in any case.
\end{quote}

\subsection*{You'll want to know: \te{సంభవించు}}
\textbf{\te{సంభవించు}} — \emph{sambhaviñcu} — \textquotedblleft{}to occur; to take place\textquotedblright{}.

\begin{quote}
\te{ఒకవేళ} \te{ప్రమాదం} \te{సంభవిస్తే}, \te{వెంటనే} \te{పోలీసులకు} \te{చెప్పాలి}.
\end{quote}

\begin{quote}
\emph{okavēḷa pramādaṁ sambhavistē, veṇṭanē pōlīsulaku ceppāli.}
\end{quote}

\begin{quote}
If an accident occurs, you must tell the police at once.
\end{quote}

Past (it …): \te{సంభవించింది}, \emph{sambhaviñcindi}.

Formal: for accidents and events, so its past is shown as \te{సంభవించింది}, it occurred. In talk, \te{జరుగు}.

\begin{cousinweb}
\textbf{\te{సంభవించు}} (\emph{sambhaviñcu}) is built on Sanskrit \emph{saṁbhava}, \textquotedblleft{}origin\textquotedblright{}, and also \textquotedblleft{}possibility\textquotedblright{} (CDIAL 12963), with the Telugu verb ending -\te{ఇంచు}.
\end{cousinweb}

\begin{grammarlens}[title={What you've built}]
One more word for saying what happens if.
\end{grammarlens}

\subsection*{Your turn}
\begin{itemize}
\raggedright
  \item \emph{Say it:} \emph{sambhaviñcu}
  \item \emph{Say it:} \emph{sambhaviñcu}, once more
  \item \emph{Say it:} \emph{okavēḷa pramādaṁ sambhavistē, veṇṭanē pōlīsulaku ceppāli.} — If an accident occurs, you must tell the police at once.
  \item \emph{Recall:} say the Telugu for an exemption, then the Telugu for in any case, then say \emph{sambhaviñcu} again
\end{itemize}

\begin{tcolorbox}[breakable,colback=teal!4,colframe=teal!35!black,title={Before you move on}]
What does \te{సంభవించు} mean? (\textquotedblleft{}To occur; to take place\textquotedblright{}.) Say it once more.
\end{tcolorbox}
